% META: {"id":"TCOMPL-CORRELATION-CAUSALITE-RE-C5","chapitre":"TCOMPL-CORRELATION-CAUSALITE","type_objet":"remediation","capacites_codes":["C5"],"status":"approved","origin":"generated","status_history":["generated","audited","accepted","approved"],"release_acceptance":"RELEASE_OWNER_BATCH_ACCEPTANCE","acceptance_closure_digest":"sha256:9b3ccf9a81c5520fbb7e03b7b2d3e7bf2057a02834b6d908bf3be5f49f3b553f","acceptance_receipt_digest":"sha256:4fad86f0282e0fa4e0419cca1ad6379ae7af5a280c04a4a90880f90a1c044242"}

%---------------------------------------------------------------
% FICHE DE REMEDIATION — C5 : identifier une variable cachee
%---------------------------------------------------------------

\begin{exercice}{TCOMPL-CORR-RE-C5-EX1}{1}{10}
Sur plusieurs pays, le nombre de téléviseurs par habitant est fortement corrélé à l'espérance de vie moyenne. Un élève en conclut que \og posséder plus de téléviseurs allonge l'espérance de vie \fg{}. Critiquer cette affirmation et proposer une variable cachée plausible.
\end{exercice}

\begin{corrige}{TCOMPL-CORR-RE-C5-EX1}
L'affirmation confond corrélation et causalité : rien ne prouve qu'un téléviseur ait un effet biologique sur la longévité. La variable cachée la plus plausible est le \textbf{niveau de développement économique} (ou richesse) du pays : un pays plus riche a à la fois davantage de téléviseurs par habitant \textbf{et} un meilleur système de santé, une meilleure alimentation, etc., qui allongent réellement l'espérance de vie. C'est cette richesse commune qui explique la corrélation observée, pas une relation directe entre télévisions et longévité.
\end{corrige}
